\section{Artifact organization and repository integration}
\label{sec:artifacts}

The ZIP is a self-contained research handoff.  Its top-level README specifies
the pinned revision, supported commands, integration paths, and the distinction
between original measurements and final audit reruns.  All performance tables
and plots in this article are generated from retained JSON rather than
transcribed by hand.  The included source files and SHA-256 manifest make
the handoff reviewable independently of the conversation that produced it.

\begin{table}[htbp]
\centering\small
\begin{tabular}{P{45mm}P{97mm}}
\toprule
Path & Purpose \\
\midrule
\path{article/} & Modular \LaTeX{} source, bibliography, compiled PDF,
generated tables, and vector/PNG figures. \\
\path{fast/} & Complete proposed Python package, with existing examples
and tests plus the new implementations and regression tests. \\
\path{baseline/fast/} & Byte-verified original files needed to execute
the archived comparison package. \\
\path{tools/} & Independent finite auditors, paired timing scripts,
table/figure regeneration, article build, and package verification. \\
\path{data/} & Raw measurements, deterministic validation records,
source provenance, final test output, and concise benchmark summaries. \\
\path{integration/} & A patch against the pinned repository paths and
notes on applying or adapting it. \\
\path{SHA256SUMS} & Hashes of packaged deliverable files; the manifest
does not hash itself. \\
\bottomrule
\end{tabular}
\caption{Research package layout.}
\end{table}

\subsection{Public behavior and compatibility}

The default recognizer now uses the checked interlacement factorizer.
The old \code{visible\_factors} function remains unchanged, and
\code{factor\_backend="legacy"} or the CLI's \code{--legacy-factor}
selects it for comparison.  New factorization evidence is a component
certificate under \code{connected\_sum\_factorization}; callers that parse
the old geometric cut trace should use the legacy backend or adopt the
new documented evidence schema.  This schema change is deliberate and
must be considered during integration.

The Alexander obstruction retains its previous witness dictionary.  Its
new backend and statistics arguments are optional.  Existing callers need
not supply them.  Dense elimination remains the reference implementation.
The runtime-budget callback now reaches sparse pivots and the factorization
stage; limits are cooperative and can still be exceeded inside an operation
without an intermediate callback, such as validation or exact polynomial work.

\code{--pointed} is an explicit recognizer option.  It currently requires
min-fill pivots, bit algebra, no deferred tail, and one scan order rather
than a race.  Incompatible CLI options receive an input error; the Python
API raises \code{ValueError}.  The direct
\code{reduced\_khovanov\_rank} result stores a \emph{reduced} rank in both
\code{rank} and \code{reduced\_rank}, and reduced values in
\code{by\_degree}.  This differs intentionally from the original
\code{khovanov\_rank} interface, whose \code{rank} is unreduced.

The separate \code{reduced\_khovanov\_decision} can return
\code{reduced\_rank=None} with a certified lower bound and replay data.
Clients must not interpret that lower bound as an exact rank.  Existing
full-rank APIs, including the legacy factored-rank helper, retain their
semantics.  The package does not silently enable pointed scanning for them.

\subsection{Reproducing the checks}

From the unpacked top-level directory, the following commands run the main
verification programs.  Explicit output names avoid overwriting the retained
research records.

\begin{lstlisting}
PYTHONPATH=fast python3 -m unittest discover -s fast/tests -v
python3 tools/verify_factorization.py \
  --output data/reproduced_factorization.json
python3 tools/verify_reduced_scan.py \
  --output data/reproduced_reduced.json
python3 tools/verify_frontier_counterexamples.py \
  --output data/reproduced_frontiers.json
python3 tools/hierarchy_budget.py --self-test \
  --output data/reproduced_hierarchy.json
python3 tools/verify_package.py
\end{lstlisting}

The recognition code and finite auditors use only the Python standard
library.  Matplotlib is needed only to regenerate figures; pre-generated
figures are included.  Building the article requires a conventional
\LaTeX{} installation with the listed packages and BibTeX, conveniently
driven by \code{latexmk}.

\begin{lstlisting}
python3 tools/render_results.py
sh tools/build_article.sh
PYTHONPATH=fast python3 -m fastunknot recognize \
  fast/examples/conway.json
PYTHONPATH=fast python3 -m fastunknot recognize \
  fast/examples/conway.json --no-alexander --no-jones --pointed
\end{lstlisting}

The benchmark README gives the timing commands, source selection, and
expected runtime considerations.  Timing scripts intentionally perform
real repeated computations and are separate from the quick regression
suite.  Reproducing exact operation counts is a stronger expectation than
reproducing identical elapsed times in a shared environment.

\subsection{Integration boundaries}

Apply the supplied patch only after checking the target revision.  If the
repository has advanced, review the differences around the touched Python
modules rather than treating the archived patch as authoritative over later
work.  The updated \path{fast/} directory is also supplied as an ordinary
source tree, so the intended changes can be inspected without applying a
patch.

Place the article and its accompanying tools/data under a new dated research
directory in \path{Topology/UnknotRecognition}, or use the repository's
preferred publication layout.  The existing synthesis remains a historical
record; the new article supplies the frontier-count correction and updated
literature status.  The bundle includes the repository's MIT-0 license.
No remote commit, pull request, or publication is required to inspect or
reproduce this handoff.
